% Responsável: TODO(grupo)
% Enunciado, item 6.2: tabela resumo padronizada, obrigatória no texto e no README.
% ATENÇÃO: nenhum número entra aqui sem vir de execução real do repositório.
\chapter{Resultados}
\label{cap:resultados}

\section{Tabela resumo}
\label{sec:tabela-resumo}

\begin{table}[htb]
  \centering
  \caption{Resumo dos resultados por técnica.}
  \label{tab:resumo}
  \footnotesize
  \begin{tabular}{llcccc}
    \toprule
    \textbf{Dataset} & \textbf{Técnica} & \textbf{Falha identificada} &
    \textbf{TPR} & \textbf{FPR} & \textbf{Acurácia/classe} \\
    \midrule
    CWRU & TODO & TODO & TODO & TODO & TODO \\
    \bottomrule
  \end{tabular}
  \fonte{Elaborado pelos autores.}
\end{table}

Técnicas que apenas detectam anomalia deixam a coluna de tipo de falha como
``não se aplica''.

\section{Detecção}
\label{sec:res-deteccao}

% Rubrica: peso 2.
TODO(grupo): curvas ROC, TPR e FPR, comparação entre os detectores one-class.

\section{Classificação}
\label{sec:res-classificacao}

% Rubrica: peso 1. Matriz de confusão COMPLETA é obrigatória (enunciado, item 6.2).
TODO(grupo). Matrizes completas no Apêndice~\ref{ap:matrizes}.

\section{Validação física}
\label{sec:res-fisica}

TODO(grupo): comparação entre o pico observado no espectro de envelope e a
frequência característica prevista para cada tipo de defeito.
